% ── CAPÍTULO 4: USO BÁSICO ───────────────────────────────────
\section{Guía de Uso Básico}
\label{sec:uso_basico}

\subsection{Tutorial en Video}

\begin{nota}[Video Tutorial]
Se ha preparado un tutorial en video que cubre la instalación, configuración y primera ejecución del benchmark. El video está disponible en el enlace que se distribuirá junto con este manual. Se recomienda verlo antes de la primera ejecución.
\end{nota}

\subsection{Formato del Archivo CSV de Entrada}

Cada red de prueba se describe mediante un archivo CSV ubicado en \lstinline|src/shared/input/|.

\begin{lstlisting}[caption={Ejemplo de estructura CSV para una red de 3 nodos.}, label={lst:csv}]
state,n0,n1,n2
0,0.8,0.3,0.6
1,0.2,0.7,0.4
2,0.5,0.5,0.9
3,0.9,0.1,0.3
...
\end{lstlisting}

\begin{itemize}[itemsep=4pt]
    \item \textbf{Columna \texttt{state}:} índice del estado presente (entero de $0$ a $2^n - 1$).
    \item \textbf{Columnas \texttt{n0}, \texttt{n1}, \ldots:} probabilidad de que el nodo $i$ tome el valor $1$ en el estado siguiente, dado el estado presente.
    \item El número de filas debe ser exactamente $2^n$.
\end{itemize}

\begin{peligro}
Si el número de filas no es una potencia de dos exacta, el sistema lanzará un error al construir el N-Cubo. Revise el conteo antes de ejecutar.
\end{peligro}

\subsection{Ejecutar el Benchmark Completo}

El script \lstinline|benchmark.py| ejecuta automáticamente todas las estrategias sobre todas las redes configuradas:

\begin{lstlisting}[style=shell, caption={Ejecutar el benchmark completo.}]
cd proyectoAlgoritmos
python src/shared/run/benchmark.py
\end{lstlisting}

El benchmark muestra el progreso en la terminal con colores:
\begin{itemize}[itemsep=2pt]
    \item Verde: estrategia ejecutada exitosamente.
    \item Amarillo: modo heurístico activado (umbral superado).
    \item Rojo: error en la ejecución (ver Sección~\ref{sec:problemas}).
\end{itemize}

\subsection{Archivos de Salida CSV}

Los resultados se guardan en \lstinline|src/shared/output/{RED}/{ESTRATEGIA}.csv|.

\begin{lstlisting}[caption={Ejemplo de archivo de salida generado.}]
index,loss,time,partition
0,0.2140,0.0082,"[[(1,0),(0,0)],[(1,1),(1,2),(0,1),(0,2)]]"
1,0.1180,0.0211,"[[(1,0)],[(1,1)],[(1,2),(0,0),(0,1),(0,2)]]"
\end{lstlisting}

\begin{itemize}[itemsep=4pt]
    \item \textbf{\texttt{index}:} número de instancia del benchmark.
    \item \textbf{\texttt{loss}:} pérdida EMD de la partición encontrada (a menor valor, mejor).
    \item \textbf{\texttt{time}:} tiempo de ejecución en segundos.
    \item \textbf{\texttt{partition}:} lista de grupos en formato Python. Cada tupla \texttt{(tiempo, índice)} representa un vértice donde \texttt{tiempo=1} es EFECTO (futuro) y \texttt{tiempo=0} es ACTUAL (presente).
\end{itemize}

\subsection{Generar el Dashboard}

\begin{lstlisting}[style=shell, caption={Generar el dashboard interactivo.}]
python src/shared/run/dashboard.py
\end{lstlisting}

Se generará el archivo \lstinline|dashboard.html| en el directorio raíz. Ábralo con cualquier navegador moderno (Chrome, Firefox, Edge).

\begin{tip}
Use \lstinline|Ctrl + Clic| en los nombres de series dentro del dashboard para mostrar u ocultar estrategias individuales y facilitar la comparación.
\end{tip}

\subsection{Interpretar la Pérdida EMD}

La columna \texttt{loss} representa cuánta información se pierde al tratar los grupos de la partición como independientes:

\begin{itemize}[itemsep=4pt]
    \item \textbf{0.0}: sin pérdida — las partes son completamente independientes.
    \item \textbf{$> 0.3$}: alta integración — el sistema tiene fuerte dependencia entre variables.
    \item La partición con \textbf{menor \texttt{loss}} es la $k$-MIP: la más informativa.
\end{itemize}
